\section*{Asesmen}
\addcontentsline{toc}{section}{Asesmen}

Asesmen Bab II bersifat formatif dan menjadi bagian dari Kuis Konsep
(3\%). Fokus asesmen adalah ketepatan notasi dan kemampuan
menghubungkan teori dengan fase kompilasi.

\begin{table}[H]
\centering
\caption{Instrumen asesmen Bab II}
\label{tab:ch02-asesmen}
\begin{tabularx}{\textwidth}{@{}lXl@{}}
\toprule
\textbf{Instrumen} & \textbf{Indikator} & \textbf{Bobot} \\
\midrule
Kuis notasi & Ketepatan penggunaan himpunan, string, dan derivasi &
formatif (bagian Kuis Konsep) \\
Tugas derivasi & Derivasi dan pohon urai benar dan lengkap &
formatif (bagian Kuis Konsep) \\
Tugas klasifikasi & Klasifikasi bahasa dan alasan tepat &
formatif (bagian Kuis Konsep) \\
Partisipasi & Keaktifan diskusi dan lokakarya & formatif \\
\bottomrule
\end{tabularx}
\end{table}

Kriteria kelulusan: mahasiswa dapat menulis definisi formal tanpa
kesalahan notasi, menghasilkan derivasi yang benar, dan menjelaskan
hubungan antara kelas bahasa dan mesin pengenal.
